% !TEX root = ../main.tex
\section{Implementation Details}
\label{app:impl}

\paragraph{Serving and decoding.}
All methods use the same vLLM~0.29 server~\citep{kwon2023vllm} and the same request settings.
Qwen3-32B is served in 16-bit precision with tensor parallelism over four RTX~3090 GPUs (24\,GB each), with prompt-lookup speculative decoding of up to 20 tokens.
Decoding is greedy with a fixed seed, at most 8{,}192 new tokens per turn, and a context window of 40{,}960 tokens.
Deterministic serving switches the server to batch-invariant mode, in which the output of a request does not depend on the requests decoded with it.
All configurations of an experiment use one server instance, and we do not rely on outputs being reproducible across server restarts.
A request that runs alone decodes at \SpdDetMsSolo\,ms per token under deterministic serving and at \SpdStdMsSolo\,ms under standard serving.

\paragraph{Cost of concurrent requests.}
We decode the first turn of four retail tasks alone and then two, three, and four of them at the same time, 320 tokens each, three times over.
For every request we compare the time it needs for the same token positions, alone and while all requests are active.
The median slowdown for two, three, and four concurrent requests is \ConDetTwo, \ConDetThree, and \ConDetFour{} under deterministic serving and \ConStdTwo, \ConStdThree, and \ConStdFour{} under standard serving.
The main runs agree with this: a parent whose successor runs during part of its turn decodes \SpdDetBusyRatio$\times$ and \SpdStdBusyRatio$\times$ slower than a turn without speculation in the two modes.

\paragraph{Speculative tool dispatch.}
The first probe is issued when the first token of a turn arrives, and further probes follow about every 50 generated tokens, at sentence boundaries, up to five per turn.
A probe may generate up to 160 tokens.
A call is dispatched when the minimum token probability over positions two through sixteen of the probe continuation is at least 0.9, and at most one call is dispatched per turn.
These settings are the same for the baseline, for the parent, and for the successor.

\paragraph{Tools and environment copies.}
The tools of \taubench run against the benchmark's local database and return within milliseconds.
All methods use the same tool latency floor within each comparison (\ToolLowFloor, \ToolMidFloor, or \ToolHighFloor\,s).
An environment copy is a deep copy of the benchmark environment.
Creating it adds a median of \MotCopyMedian\,s between the dispatch decision and the start of the tool.
The fingerprint of an environment is the pair of hashes of its agent-side and user-side databases.

\paragraph{The check.}
Observations are compared by SHA-256.
The configuration compared by condition~(v) consists of the model, the endpoint, the seed, the token budget, the context window, the thinking flag, the stop strings, the temperature, and the presence penalty.

\paragraph{Successor.}
The successor uses the settings of a regular turn with one exception.
It does not reuse probe tokens as drafts, because our server keeps a single draft state that is shared by all requests.
The parent keeps this optimization.
The drafts that the successors forgo would save at most \MotDraftForgonePct\% of the time.

\paragraph{History.}
For all methods, the assistant message stored in the history is the output of the turn with its reasoning block removed.
